\documentclass[11pt]{article}
\usepackage[letterpaper,margin=.86in]{geometry}
\usepackage[T1]{fontenc}
\usepackage{lmodern,microtype,amsmath,amssymb,amsthm,mathtools,booktabs,longtable}
\usepackage[dvipsnames]{xcolor}
\usepackage{enumitem,fancyhdr,xurl,hyperref}
\hypersetup{colorlinks=true,linkcolor=MidnightBlue,citecolor=MidnightBlue,urlcolor=MidnightBlue,pdftitle={Two reductions for Diophantine consistency},pdfauthor={AI-assisted research prepared for Vladimir Reshetnikov}}
\setlength{\parindent}{0pt}
\setlength{\parskip}{5pt}
\setlength{\emergencystretch}{2em}
\setlist{nosep,leftmargin=*}
\allowdisplaybreaks
\newtheorem{theorem}{Theorem}[section]
\newtheorem{lemma}[theorem]{Lemma}
\newtheorem{proposition}[theorem]{Proposition}
\theoremstyle{remark}\newtheorem{remark}[theorem]{Remark}
\newcommand{\Z}{\mathbb Z}
\newcommand{\N}{\mathbb N}
\newcommand{\Zp}{\mathbb Z_{>0}}
\newcommand{\Con}{\operatorname{Con}}
\newcommand{\fl}[1]{\left\lfloor#1\right\rfloor}
\newcommand{\rg}[2]{\{#1,\ldots,#2\}}
\newcommand{\pin}{\texttt{84e16e43092614bae262ca59651cb9d77e31a255}}
\newcommand{\file}[1]{{\small\nolinkurl{#1}}}
\pagestyle{fancy}\fancyhf{}
\fancyhead[L]{\small Two reductions for Diophantine consistency}
\fancyhead[R]{\small 24 October 2026}
\fancyfoot[C]{\thepage}\setlength{\headheight}{14pt}
\title{\vspace{-1.5em}\textbf{Two reductions for Diophantine consistency}\\[.4em]
\large A two-root equation of degree 278 and a smaller direct integer quartic}
\author{AI-assisted research prepared for Vladimir Reshetnikov\\
\small Integration packet for the ProveIt research program}
\date{24 October 2026}

\begin{document}
\maketitle\thispagestyle{empty}
\begin{abstract}
We give two explicit reductions relative to the ordinary calibrated source
theorems in a fixed ProveIt snapshot. The first gives an integer-coefficient
polynomial in ten strictly positive unknowns with exact degree 278 and a
93-operation evaluator. Its immediate predecessor has degree 368 and 99
operations. The construction saturates a strong Pell discriminant by the first
Pell radicand, recovers two integral roots from their coprime product, and uses
an upper polynomial center to shorten the transport condition. Two exact
center identities establish the degree.

The second changes the shared numerical bound in the direct proof-certificate
construction. Squaring selected signed outputs permits recovery of their
intended powers before imposing their signs. Together with additional exact
source consequences this yields a quartic in 576 unrestricted integers with
7,172 terms; a smaller-coefficient companion has 582 integers and 6,200 terms.
These improve the corresponding 604- and 609-variable source equations.
Every coefficient, residual, arithmetic circuit and reconstruction certificate
is supplied in the packet.

Both reductions preserve existence of a finite ZFC refutation. Composing the
ordinary source calibrations gives ZFC-provable equivalences between
unsolvability and syntactic $\Con(\mathrm{ZFC})$. The new algebra is checked
exactly, but no new end-to-end Lean/Rocq or encoded object-ZFC proof is claimed.
Neither global optimality nor the consistency of ZFC is asserted.
\end{abstract}

\newpage
\tableofcontents
\newpage

\section{Results, baseline, and meaning of simplicity}

All source comparisons use ProveIt commit \pin. This is essential: the active
research is considerably ahead of the numbered manuscript collection. The
99-operation zero-offset equation and the 75-predicate direct quartic family
are the immediate baselines, not the older 101-operation printed page or
earlier hundreds-of-operation norm constructions \cite{zerooffset,output75}.

\begin{center}\small
\begin{tabular}{@{}llrrr@{}}
\toprule
Construction & Domain & Unknowns & Degree & Operations\\
\midrule
Zero-offset parent & $\Zp$ & 10 & 368 & 99\\
Two-root equation, this packet & $\Zp$ & 10 & 278 & 93\\
Simple integer substitution & $\Z$ & 30 & 556 & 163\\
\bottomrule
\end{tabular}
\end{center}

The first improvement saves six live arithmetic operations and lowers degree
by 90, or about $24.5\%$, without adding a positive unknown. Operations are
binary additions, subtractions and multiplications, including squaring and
multiplication by a small numeral. Named polynomial abbreviations are not
additional unknowns. Fixed coefficient construction is accounted for
separately.

\begin{center}\small
\begin{tabular}{@{}lrrr@{}}
\toprule
Unrestricted integer quartic & Unknowns & Terms & Coefficient height\\
\midrule
Matched prior small-variable profile & 604 & 9,676 & 2,445,790,974,704\\
New small-variable profile & 576 & 7,172 & 2,445,790,974,704\\
Matched prior smaller-height profile & 609 & 8,823 & 964,213\\
New smaller-height profile & 582 & 6,200 & 964,213\\
\bottomrule
\end{tabular}
\end{center}

Here a term means a nonzero collected monomial, and height is the maximum
absolute coefficient of the complete expanded polynomial. Thus the first
quartic saves 28 variables and 2,504 terms at the same height. The second saves
27 variables and 2,623 terms at the same height. They are two measured
tradeoffs, not a claim to have characterized a complete Pareto frontier.

There is no established scalar ordering that makes all these objects
``simplest.'' A quartic in many integer variables and a degree-278 polynomial
in ten positive variables answer different questions. The repository also
retains fewer-variable integer constructions based on Sun's reduction, and
nine-positive constructions with much larger degree. Sun's general
ten-integer theorem likewise is not displaced by a ten-positive result
\cite{sun2026}. The present contributions are concrete improvements on two
specified routes.

\subsection{What is proved inside ZFC}

Fix the syntactic ZFC proof predicate used by the repository and write
$\Con(\mathrm{ZFC})$ for the assertion that it has no finite refutation.
The desired statement for an integer polynomial $P$ and a domain $D$ is
\[
\mathrm{ZFC}\vdash
\bigl(\forall\mathbf{x}\in D^m\;P(\mathbf{x})\ne0\bigr)
\longleftrightarrow\Con(\mathrm{ZFC}).
\]
This is an equivalence, not a proof of either side. It does not contradict
the second incompleteness theorem. We give ordinary mathematical derivations
of the new arithmetic equivalences and explicitly name the inherited source
calibrations. The finite algorithms used in the proofs, including recurrence
periods, are justified by ordinary finite arithmetic and finite group
arguments within ZFC; no standard-model-only soundness premise is added.

Three evidence levels should remain distinct. The new polynomial identities
and all delivered coefficients have exact executable checks. The unbounded
soundness and completeness arguments are ordinary proofs. The repository's
native/calibration and internal-ZFC formalization boundaries remain as
documented in the source snapshot. This packet does not close those formal
development tasks merely by running a polynomial checker.

\section{The complete ten-positive-variable equation}
\label{sec:equation}

The only unknowns are
\[
J,F,\alpha,i,T,w,Y,n,\eta,\zeta\in\Zp.
\]
The four fixed integers are $\beta,K_0,M_C,M_F$. Appendix~\ref{app:constants}
gives a closed finite recipe, including every seed character. All symbols
below are polynomial abbreviations. Put
\begin{align*}
q&=\beta J+1,&C&=q-F-\alpha,&X&=wq,\\
k&=\eta+\zeta,&c&=kY+\eta,&\mu&=q^2+1,\\
A&=\mu Y(X+1)+1,&\Delta&=A^2-1,&H&=4A-5,\\
V_0&=X\mu^2Y^2,&\ell&=V_0k,&D&=\ell(\ell+k)+1.
\end{align*}
The source index, transport expression and index error are
\begin{align*}
R&=[(q-1)(q-F)+\alpha](q^2-1)+(M_C+qM_F)J,\\
B_t&=(K_0+w)C-F,& e&=k-R-1.
\end{align*}
The changed strong and auxiliary quantities are
\begin{align*}
I_0&=ic^2,& B_f&=1+\Delta D I_0^2,&Q&=\Delta^2 D I_0^2,\\
\delta&=Q-1,&b&=c(T+\alpha)-R,&V&=B_fb-c,\\
C_y&=QV^2-1,&P&=AI_0D,&W&=1+(q-1)b.
\end{align*}
Finally define
\begin{align*}
K_+&=X(3c-2X)+2, &P_1&=q^3(q-1), &M&=P_1\mu^2H,\\
Z&=P_1\mu^2[H(n+V+P)+K_+]\\
 &\hspace{1.4em}+H[\mu^2((q-1)Y-q^3WB_t)-P_1e],\\
s&=\delta Z^2,&u&=\delta M^2DB_f,&v&=M^2C_y.
\end{align*}
The requested single equation on this domain is
\begin{equation}\label{eq:main}
\boxed{\mathcal P=(s-u-v)^2-4uv=0.}
\end{equation}

In particular, this is an ordinary polynomial equation with fixed integer
coefficients. The radicals and divisions used in its proof do not occur in
its evaluation. The delivered arithmetic graph has 55 multiplications,
22 additions and 16 subtractions; every row, all ten inputs and all four
coefficient ports are live.

The whole BFS-renamed compact44 coefficient tuple is used. It satisfies
\begin{equation}\label{eq:portbox}
\beta\text{ odd},\quad\beta\ge31,\quad K_0>0,
\qquad 0<M_C<\beta<M_F<2\beta.
\end{equation}
The values of $\beta,M_C,M_F$ agree with the original compact44 tuple, but
$K_0$ differs under the coherent state/color renaming. We use its separate
source transfer, not a numerical identification of the two $K_0$ values
\cite{bfspage}. The printed recipe has
$\beta=2^{590069}-1$; the four payload bit lengths are
$590069,592884,590069,590070$ for $\beta,K_0,M_C,M_F$, respectively.
The small operation count should not conceal these coefficient sizes.

\begin{theorem}\label{thm:ten}
For the complete fixed tuple in Appendix~\ref{app:constants},
Equation~\eqref{eq:main} has a solution in $\Zp^{10}$ if and only if
the calibrated source-61 system has a positive solution. By its inherited
ordinary calibration, this is equivalent to the existence of a finite
syntactic ZFC refutation. The exact total degree of $\mathcal P$ is 278.
\end{theorem}

\section{Two elementary extraction lemmas}

\begin{lemma}[Unequal two-root extraction]\label{lem:two}
Let $r,y>0$ have rational squares. If $r\ne y$ and a signed sum of
$r,y$ is rational, then both $r$ and $y$ are rational.
\end{lemma}
\begin{proof}
Write $x=\varepsilon r+\eta y\in\mathbb Q$, with both signs in
$\{1,-1\}$. Since $r\ne y$, $x\ne0$. Then
$2\varepsilon xr=x^2+r^2-y^2$, so $r$ is rational, and the equation
for $x$ gives rationality of $y$.
\end{proof}

\begin{lemma}[Adjacent-coefficient descent]\label{lem:adjacent}
Let $Q>1$ and $V$ be integers. If $y\in\mathbb Q$ satisfies
$QV^2-(Q-1)y^2=1$, then $y$ is an integer.
\end{lemma}
\begin{proof}
Replace $V,y$ by their absolute values. Neither is zero. Write $y=a/m$
in lowest terms with $m>0$. The equation implies $m^2\mid Q-1$.
For $V=1$ we have $y=1$. For $V>1$, set
\[
V'=(2Q-1)V-2(Q-1)y,\qquad
y'=(2Q-1)y-2QV.
\]
The number $V'$ is an integer and the same norm equation holds for
$(V',y')$. Moreover $y>V$ and $QV>(Q-1)y$, by subtracting the relevant
squares. Consequently $-V<V'<V$, and the norm excludes $V'=0$.
The reduced denominator of $y'$ is still $m$, because
$2Q-1\equiv1\pmod m$. Strong induction on $|V|$ reaches $V=1$ and
forces $m=1$.
\end{proof}

Square-root saturation itself has a predecessor in the repository's
nine-positive route \cite{ninesaturation}. The following use of it avoids
the weighted three-root noncancellation argument of the ten-positive parent.

\begin{proposition}[The saturated certificate]\label{prop:saturation}
Every positive zero of \eqref{eq:main} yields positive integers $\tau,f,y$
such that
\[
D=\tau^2,\qquad B_f=f^2,\qquad
f^2-\Delta(I_0\tau)^2=1,\qquad
QV^2-(Q-1)y^2=1.
\]
In particular, $I=I_0\tau$ is a genuine strong Pell ordinate and
$c^2\mid I$.
\end{proposition}
\begin{proof}
At every positive input, $q\ge32$, $A>4$, $c\ge3$, $I_0\ge c^2$ and
$D>1$. Directly from the definitions,
\[
Q-1>c^2-1>0,\qquad B_f>c+2,\qquad B_f<Q-1.
\]
If $b=0$, then $V=-c$; otherwise $|V|\ge B_f-c>1$. Thus $C_y>0$.
Set $r=\sqrt{DB_f}$ and $y=\sqrt{C_y/\delta}$ as positive real roots.
The four-conjugate identity for $(x^2-a-b)^2-4ab$ gives
\begin{equation}\label{eq:rootrelation}
Z/M=\varepsilon r+\eta y.
\end{equation}
Its prefactor $\delta^2M^4$ is positive. If $r=y$, clearing denominators
would give $(Q-1)DB_f=QV^2-1$. Reduction modulo $D$ gives $0=-1$,
because $D\mid Q$, a contradiction.

Lemma~\ref{lem:two} gives rational $r,y$. A rational number with integer
square is an integer, so $r\in\Z$. Since $B_f\equiv1\pmod D$,
$\gcd(D,B_f)=1$. Prime valuations in the integer square $r^2=DB_f$
show that both factors are squares. Lemma~\ref{lem:adjacent} gives
$y\in\Z$. Only now define $I=I_0\tau$; its stated properties follow.
\end{proof}

\section{Soundness: restoring the source in the right order}

For $A\ge2$, write
\[
\chi_A(j)+\psi_A(j)\sqrt{A^2-1}
       =(A+\sqrt{A^2-1})^j.
\]
The inherited normalized source proof supplies the strong-rank and signed
auxiliary-index lemmas, and the immediate zero-offset parent supplies its
ratio and projection refinements \cite{source61,zerooffset}. Their exact
hypotheses, rather than just their conclusions, are used below.

\subsection{Folding and first growth before transport}

The strong equation and auxiliary residues
\[
V\equiv-c\pmod f,\qquad c\le I,
\]
give $c=\psi_A(p)$ for a positive odd index $p$. One way to see the folding
step is to reduce an auxiliary Pell index modulo the strong period:
$\psi_A(j+2m)\equiv-\psi_A(j)\pmod{\chi_A(m)}$ and
$\psi_A(2m-j)\equiv\psi_A(j)$. The representatives have size at most $I$,
whereas $f>2I$. Thus the residue of $c$ determines a positive representative.
The odd auxiliary index preserves its parity through this folding.

The first norm gives
$k=2\psi_{2V_0+1}(n_0)$ for $n_0\ge1$. If $n_0=1$, then $k=2$,
$\eta=\zeta=1$, and $c=2Y+1$ lies strictly between $\psi_A(1)$ and
$\psi_A(2)$, which is impossible. Hence
\begin{equation}\label{eq:earlyc}
c>kY\ge k\ge4(2V_0+1)>8X\mu^2Y^2.
\end{equation}
No transport conclusion, parity conclusion for $q$, or divisibility of $e$
has yet been used. In particular, $c>2X$ now proves $K_+>0$.

The unchanged mask estimate is
$R<q^4+(\alpha+2)q^2$. Equation~\eqref{eq:earlyc} gives
$R<c\alpha/4$ and hence $b>0$. The same two cases $R\ge0$ and $R<0$
as in the parent give $|e|<b$, so $|e/\mu^2|<b/2$.
Since $B_f<Q-1$ and $V>0$,
\[
0<y-V=\frac{V^2-1}{(Q-1)(y+V)}<\frac b2.
\]

\subsection{The upper center removes a large transport term}

The new center has the exact defect
\begin{equation}\label{eq:defect}
P^2-r^2=D(I_0^2D-1)>0.
\end{equation}
Thus $r<P$. Dividing the definition of $Z$ by $M$ in the proof gives
\[
Z/M=n+V+P+K_+/H+Y/q^3-WB_t/(q-1)-e/\mu^2.
\]
Put $t=B_t/(q-1)$ and $\rho=e/\mu^2$ temporarily as rationals. The
signed-root equation implies
\begin{equation}\label{eq:ntransport}
n=\varepsilon r-P+(\eta y-V)-K_+/H-Y/q^3+Wt+\rho.
\end{equation}
If $B_t\le-1$, then $Wt<-b$. The terms $\varepsilon r-P$,
$-K_+/H$ and $-Y/q^3$ are negative; the other two errors together are
less than $b$. Equation~\eqref{eq:ntransport} would give $n<0$.
Consequently $B_t\ge0$, and since $F>0$ and $K_0+w>0$, also $C>0$.
The source masks now give
\begin{equation}\label{eq:Rbounds}
3q<R<q^4-2.
\end{equation}

This is the important use of the upper center. The old weight contained
a large strong-ordinate term. Here $W=1+(q-1)b$ suffices. Its degree is 5,
which will be used in the exact degree calculation.

\subsection{Rank, parity, projection, and the unscaled source}

For reference, the normalized rank argument is short. If
$c=\psi_A(p)>2$ and $c^2\mid\psi_A(m)$, strong divisibility gives
$m=pv$. Put $D_m=\chi_A(p)$. Then
$\psi_A(m)=c\psi_{D_m}(v)$ and
$\psi_{D_m}(v)\equiv vD_m^{v-1}\pmod c$, with $D_m$ a unit modulo $c$.
Thus $c\mid v$, so $pc\mid m$ and $m>2p$.

Apply this to the recovered $I=I_0\tau$. The signed auxiliary-index
lemma also uses $V\equiv-R\pmod c$, which follows from $f^2\equiv1\pmod c$.
Its strict bounds $0<p,R<c/2$ follow from \eqref{eq:earlyc},
\eqref{eq:Rbounds}, and Pell growth. It therefore gives $p=R$.
Since $p$ is odd, $R$ is odd. Odd $q$ would make $J$ even and the defining
expression for $R$ even, because $\beta$ is odd. Thus $q$ is even.

Only at this point use
$\gcd(\mu^2,q^3(q-1)H)=1$. In the signed-center identity,
$n,V,P,r,y$ are integers, so a partial-fraction argument yields
$\mu^2\mid e$. The first-Pell congruence and the comparison
$2V_0+1>A$ now force $n_0=(R+1)/2$. Write $R=2h+1$; then $h\ge48$.
For $a=Y(X+1)$, the zero-offset ratio estimate is
\[
\xi<\frac{\psi_A(2h+1)}{\psi_{2V_0+1}(h+1)}
<\xi(1+1/a)^{2h},\qquad
\xi=\frac{(X+1)^{2h}}{X^h}.
\]
It gives $Y>X^h/3$ and $X^{h+1}<H$.

Set $K_{\rm quad}=X(X+2(A-2)c)-1$. The exact identity
$K_+=HXc-2K_{\rm quad}$ and the oddness of $H$ transfer the projection
divisibility to $K_{\rm quad}$. The Pell projection then yields
\[
H\mid q^3(q-1)(X-2^R)(2^RX+1).
\]
For $X\ge q\ge32$ and $h\ge48$, its entire absolute value is less than
\[
q^4 2^{2R+2}X^2<X^{h+1}<H.
\]
The dividend is therefore zero, and its second factor is positive, so
$X=2^R$. No nonunit factor is cancelled through the composite modulus.
The remaining partial fractions give $q^3\mid Y$ and $q-1\mid B_t$.
The retained fixed-minus parity lemma gives $R\equiv3\pmod4$, and ratio
rounding recovers the required half-binomial coefficient $Y$.

These are exactly the data from which the parent constructs fresh unscaled
source-61 witnesses. The reconstruction does not divide the scaled base or
reuse the current strong witness as an unscaled one. Its calibrated source
theorem, including the allowed positive $J=1$ case, yields a finite
refutation. This proves soundness.

\section{Completeness: the new witnesses really exist}

It is not enough to retain the predecessor's witnesses: the new strong
ordinate must be divisible by $\tau$, and the upper center subtracts a large
positive term from the expression for $n$. Both points need proof.

Canonical source data give $X=2^R$, $R\equiv3\pmod4$,
$q^3\mid Y$, and $B_t=(q-1)t$ with $t\ge0$. Retain $J,F,\alpha,w,Y$,
and choose
\[
c=\psi_A(R),\quad
k=2\psi_{2V_0+1}((R+1)/2),\quad
\tau=\chi_{2V_0+1}((R+1)/2).
\]
The unchanged ratio argument gives $Y<c/k<Y+1$, so
$\eta=c-kY$ and $\zeta=k-\eta$ are positive integers. Also
$\rho=(k-R-1)/\mu^2$ is a positive integer.

\begin{lemma}[Divisible strong witnesses]
For any positive integer $L$ and $A\ge2$, there is a positive index $m$
with $L\mid\psi_A(m)$.
\end{lemma}
\begin{proof}
The matrix $\begin{pmatrix}A&A^2-1\\1&A\end{pmatrix}$ has determinant one.
Its reduction modulo $L$ belongs to a finite group, so some positive power
is the identity. The lower-left entry of its $m$th power is $\psi_A(m)$.
\end{proof}

Apply the lemma with $L=c^2\tau$ and set
\[
I=\psi_A(m),\qquad f=\chi_A(m),\qquad
i=I/(c^2\tau),\qquad I_0=I/\tau.
\]
These are positive integers, $I_0=ic^2$, and the saturated strong equation
holds. No index variable or exponent is introduced into the polynomial.

Put $S=\Delta I$ and $Q=S^2$. The standard odd-index auxiliary witnesses are
\[
V_*=\frac{\chi_S(R)}{S},\qquad y_*={\psi_S(R)},\qquad
T_*=\frac{V_*+c+Rf^2}{cf^2}.
\]
The retained odd-index polynomial identities give
$V_*\equiv-c\pmod{f^2}$ and $V_*\equiv-R\pmod c$.
Since $f^2\equiv1\pmod c$, $T_*$ is an integer and
$V_*=f^2(cT_*-R)-c$.

Here $R\ge5$. Monotonicity and the explicit fourth and fifth Pell terms give
\begin{align*}
V_*&\ge\chi_S(5)/S=16Q^2-20Q+5>Q^2,\\
u_*:=y_*-V_*&=\frac{\psi_S(R-1)}{S}\ge8Q-4.
\end{align*}
The exact gap identity follows from
$\chi_S(R)-S\psi_S(R)=-\psi_S(R-1)$.
Since $f^2<Q$ and $Q/c>q>\alpha$, $T_*>\alpha$.
Thus $T=T_*-\alpha$ is a positive integer and the equation's $V$ equals $V_*$.

The projection gives the positive integer $\kappa=K_+/H<Xc$.
Set $s_0=Y/q^3\ge1$ and define
\begin{equation}\label{eq:complete_n}
n=Wt+\rho+\tau f+u_*-P-s_0-\kappa.
\end{equation}
It is an integer. Moreover $P<Q$ because
$Q/P=\Delta^2I_0/A>1$. Equation~\eqref{eq:earlyc} implies
$X<c/8$ and $Y<c$, hence
$s_0+\kappa<Xc+Y<c^2<Q$. Therefore
\[
P+s_0+\kappa<2Q<8Q-4\le u_*.
\]
All other terms in \eqref{eq:complete_n} are nonnegative, so $n>0$.
The normalized center now equals $\tau f+y_*$, making one factor of
\eqref{eq:main} vanish. This proves completeness, including positivity of
both changed witnesses $i$ and $n$.

\section{Exact degree and an unrestricted-integer version}

The direct syntactic degree propagation of the 93-gate graph gives 316.
The exact degree is smaller. Introduce proof-only polynomial abbreviations
\[
z=Z-M(V+P),\qquad d=P^2-DB_f=(I_0D)^2-D,
\]
and
\[
B=2M(P+V)z+z^2+M^2d,\qquad
C_0=\delta B-M^2(V^2-1).
\]
A complete ring identity is
\begin{align}\label{eq:centeridentity}
\mathcal P={}&C_0^2+4\delta M^2PV C_0
-4\delta M^4P^2(V^2-1)\nonumber\\
&+4\delta M^4d(QV^2-1).
\end{align}
It follows by substituting $Z=M(V+P)+z$ and $DB_f=P^2-d$ into the
norm and using $\delta=Q-1$; the checker compares every coefficient of
this six-indeterminate identity.

The relevant degrees are
\[
\deg Q=48,\quad\deg M=13,\quad\deg V=42,\quad\deg P=28,
\quad\deg z\le19,\quad\deg d=46.
\]
Thus $\deg B\le74$ and $\deg C_0\le122$. The four terms in
\eqref{eq:centeridentity} have degrees at most $244,266,240,278$.
The only degree-278 contribution is $4Q_0^2M_0^4d_0V_0^2$, where subscripts
denote leading homogeneous forms. To avoid confusing $V_0$ here with the
earlier register, its fully expanded factorization is
\begin{equation}\label{eq:leading}
2^{10}q_0^{146}i^{10}w^{28}Y^{60}(\eta+\zeta)^{32}(q_0-F)^2,
\qquad q_0=\beta J.
\end{equation}
In particular the coefficient of
$J^{146}i^{10}w^{28}Y^{60}\eta^{32}F^2$ is
$2^{10}\beta^{146}>0$. The exponents sum to 278, proving the claimed
exact degree. This proof is independent of the numerical values of
$K_0,M_C,M_F$.

For an immediate unrestricted-integer equation, substitute, for each of
the ten positive inputs,
\begin{equation}\label{eq:legendre}
x_j=1+a_j^2+b_j^2+c_j(c_j+1),\qquad a_j,b_j,c_j\in\Z.
\end{equation}
This expression parametrizes precisely the positive integers. The converse
uses the three-square theorem on $4(x_j-1)+1$, with one odd and two even
squares. Thus the resulting 30-variable equation has the same refutation
equivalence. Its degree is exactly 556: substitution of independent
nonzero quadratic leading forms into the nonzero product
\eqref{eq:leading} stays nonzero. The stated evaluator adds seven operations
per input, for $93+70=163$ operations. This companion is explicit and
useful for domain clarity, but it is not advertised as a new unrestricted
variable-count or low-degree record.

\section{A direct consistency quartic in 576 unrestricted integers}
\label{sec:direct576}

The direct proof-certificate route admits a new simplification: a
quadratic bound on a proposed power output can replace its sign
predicate.  Once the indexed Pell coordinates are oriented, a small
signed residue identifies the output and recovers its sign.  Combining
this with two local source identities and a factorial digit obstruction
gives the following explicit quartics.

\begin{center}\small
\begin{tabular}{lrrrr}
\hline
Equation & Variables & Quadratic rows & Terms & Height\\
\hline
Previous principal &604&372&9,676&2,445,790,974,704\\
New principal &\textbf{576}&368&\textbf{7,172}&2,445,790,974,704\\
Previous lower height &609&377&8,823&964,213\\
New lower height &\textbf{582}&374&\textbf{6,200}&964,213\\
\hline
\end{tabular}
\end{center}

All four total degrees are exactly four.  The two new equations are
literally
\[
 P_{576}(x)=\sum_{j=0}^{367}R_j(x)^2=0,
 \qquad
 P_{582}(x)=\sum_{j=0}^{373}\widetilde R_j(x)^2=0,
\]
where the accompanying JSON files list every quadratic residual and
every nonzero expanded coefficient.  There are no coefficient ports or
unmaterialized variable powers.  The first equation saves 28 variables
and 2,504 monomials at unchanged height; the second saves 27 variables
and 2,623 monomials.  These are measured improvements, not a minimality
claim.  This direct C222 route does not use K215, Jones compression or
Sun's integer reduction.

The precise new theorem is integer-zero existence equivalence with the
pinned 75-predicate direct-C222 source.  Its existing ordinary
refutation/consistency calibration transfers separately.  All arguments
here are finite integer arithmetic that can be performed in ZFC; no
consistency or external soundness assumption is made.  A new Lean,
Rocq or sealed object-language ZFC proof object is not claimed.

\subsection{The exact source change}

The source is pinned at ProveIt commit
\texttt{84e16e43092614bae262ca59651cb9d77e31a255}, in the
\texttt{factorial-radius-output-sign-profile/raw} research packet.
Its body has 1,259 coordinates and 1,177 quadratic residuals.  A
predicate $E\ge0$ is represented by
\[
 E-a^2-b^2-c(c+1)=0
\]
with three private integer variables.  These witnesses exist exactly
when $E\ge0$, by the three-square theorem applied to $4E+1$.

Delete exactly the twelve scalar predicates with raw indices
\[
 6,104,106,575,584,592,600,609,618,627,636,825.
\]
Keep the other 63 predicates.  Put
\[
 S=\{6,106,573,583,591,599,607,617,625,635\}.
\]
Change only body residual 1176, originally $x_{643}-x_{1258}$, to
\[
 x_{643}-x_{1258}+\sum_{i\in S}(x_i-x_i^2).
\]
This adds no coordinate and preserves quadratic degree.  Every other
body residual is unchanged.  The resulting raw equation has 1,448
integer coordinates and 1,240 quadratic residuals.

Let $\mathcal A$ be the set of all 43 distinct base, exponent and
output indices in the 21 actual power calls.  The original frame
gate identities now give exactly
\begin{equation}
 W=2+\sum_{i\in\mathcal A\setminus S}x_i+\sum_{i\in S}x_i^2,
 \qquad W=x_{643}.
 \label{eq:direct-signed-frame}
\end{equation}
The complete list $\mathcal A$ and the actual source-row ideal are
included in the certificate.  A squared contribution bounds either
sign of an integer because $z^2\ge |z|$.

\subsection{Signs and a dominating Pell frame}

Write $f=x_{87}$, $r=x_{94}$, $s=x_{95}$ and $N=x_{93}$.  The
retained predicates give $f\ge2$, $r,s\ge0$.  Exact source rows
78--85 imply
\[
 N=fr+(f-1)s+f\ge2.
\]
Separately, an eleven-row polynomial identity gives
\begin{equation}
 f(x_{104}+1)=f x_{59}+x_{131}+3.
 \label{eq:direct-output104}
\end{equation}
The two right-hand coordinates retain direct nonnegativity predicates.
Thus $x_{104}+1>0$ and $x_{104}\ge0$, without any power interpretation.
The full eleven integral polynomial multipliers are supplied.

Filtering the old prefix proof now proves 604 of the 643 semantic
coordinate signs.  Its only postponed special bound concerns $x_{130}$.
The only unknown-sign frame arguments are precisely the ten members
of $S$.  All actual bases and seed exponents are already positive.
In particular, with $t=x_4$, $K=x_{86}$ and $f_1=x_{107}$,
\begin{align*}
 &t\ge1,\quad K\ge3,\quad f_1\ge0,\\
 &M=x_{129}=1+f(x_{59}+1)\ge3,\\
 &b=x_{530}=M(1+N+x_{113})\ge9,\qquad x_{531}=b+1\ge10,\\
 &x_{103}=1+(t+1)f\ge5.
\end{align*}
Consequently the new frame gives, for every actual call $(n,k,m)$,
\begin{equation}
 W\ge3,\qquad 1\le n,k<W,\qquad |m|<W.
 \label{eq:direct-frame-bounds}
\end{equation}

The retained initial modulus and its identity
$T_0=2A(t+1)-(t+1)^2-1$ force $A=x_{644}\ge2$.
The unchanged bootstrap norm is
\[
 A^2-1=W^3(W+2)Z^2.
\]
Here $Z\ne0$.  At Pell base $W+1$ this is the solution
$(A,W|Z|)$.  Classification and the recurrence modulo $W$ imply
that its positive index is a multiple of $W$.  Hence
\begin{equation}
 A>W^W,
 \qquad T_n=2An-n^2-1\ge A>W^W.
 \label{eq:direct-large-modulus}
\end{equation}
The last inequality follows from the exact decomposition
\[
 T_n-A=(2n-1)(A-n-1)+(n-1)(n+2).
\]

The retained $V\ge1$ predicate and the unchanged product identity
$V=(\prod_iY_i^2)(v+1)$ make every seed ordinate nonzero and give
$Y_i^2\mid V$.  The common norm implies $|U|\ge A$.  From
$B=1+4V(q+1)$ and $B\equiv A\pmod U$, the case $q=-1$ is impossible:
it would force $U\mid A-1$, although $0<A-1<|U|$.  Every other $q$
gives $|B|\ge3$.  Thus the pinned signed-base index theorem is
available whenever the particular short-interval condition
$k\le |Y|$ has been established.

\subsection{A power output may initially be signed}

\paragraph{Signed-output lemma.}
In this section $\chi_k(a),\psi_k(a)$ use an index subscript:
they start at $(1,0)$ and satisfy
$\chi_{k+1}=a\chi_k+(a^2-1)\psi_k$,
$\psi_{k+1}=\chi_k+a\psi_k$.
Assume \eqref{eq:direct-frame-bounds}--\eqref{eq:direct-large-modulus},
$n\ge2$, and
\[
 |X|=\chi_k(A),\quad Y=\psi_k(A)>0,\qquad
 X-(A-n)Y-m=T_nh,
\]
where $m,h$ are unrestricted integers.  Then
$X=\chi_k(A)$ and $m=n^k$.

\paragraph{Proof.}
Put $p=\chi_k-(A-n)\psi_k$ and
$q=\chi_k+(A-n)\psi_k$.  The Pell recurrence and norm give
$p\equiv n^k\pmod{T_n}$ and $pq=1+T_n\psi_k^2$.
If $X<0$, then $T_n\mid mn^k+1$.  This dividend is nonzero
for every integer $m$, since $n^k\ge2$, but
\[
 |mn^k+1|\le(W-1)^W+1<W^W<T_n.
\]
This is impossible.  Hence $X>0$, and $T_n\mid m-n^k$.  Now
\[
 |m-n^k|\le W-1+(W-1)^{W-1}<W^W<T_n,
\]
so $m=n^k$.  Notice that positivity of $Y$ is essential:
if both main coordinates could be negative, $m=-n^k$ would survive.
Every application below first supplies that positive ordinate.

\subsection{Restoring the deleted predicates in order}

\paragraph{The eight range quotients.}
For actual calls $10,12,13,14,15,17,18,20$, all eight original
$Y\ge k$ predicates remain.  The respective outputs are
$573,583,591,599,607,617,625,635$; all bases are $x_{531}\ge10$.
The signed index theorem followed by the signed-output lemma makes
these outputs their intended positive powers.  Their exact range
rows are
\[
 m=(Me)q+r,\qquad e=r+s+1,\qquad r,s\ge0.
\]
Here $M=x_{129}$, not the radix $b=x_{530}$.  Since
$0\le r<e\le Me$, a negative $q$ would make $m<0$.
All eight quotients and their 24 dependent signs are recovered.
Only the seven prefix signs
\[
 6,106,124,125,130,132,133
\]
remain unproved.

\paragraph{The first factorial.}
None of those seven coordinates is a sign premise of the pinned
local first-factorial argument.  Put $d=fK^2$ and $C=x_{91}$.
Its exact inputs are
\[
 C=r+s+1,\quad fC=N+s,\qquad
 x_{532}=x_{533}(C+ba)+u,
\]
\[
 x_{533}=u+v+1,\qquad b=C+z+1,
\]
with $a,u,v,z\ge0$, together with actual calls $1,5,6$.
All three outputs $N,x_{532},x_{533}$ remain known nonnegative.

For either missing first-factorial index, $|Y|\ge2$ implies
$|Y|\ge2A>W>k$, so the index theorem and a nonnegative-output
supporting call recover its true positive pair.  If $|Y|=1$, the
four possible residues $n,-2A+n,2A-n,-n$ and the small-output bound
leave exactly $X=A,Y=1,m=n$.  Thus the effective indices are
$e\in\{K,1\}$ and $j\in\{K+d,1\}$, with
\[
 N=(d+1)^e,\quad x_{533}=b^e,\quad x_{532}=(b+1)^j.
\]
If $j=1,e=K$, the digit equation contradicts $C\ge1$; if
$j=e=1$, it forces $C=1$ and then $f\ge d+1>f$.
Therefore $j=K+d$.  If $e=1$, the bound
$b\ge(f+1)(d+2)>K+d$ identifies the first binomial digit as
$C=K+d$, and the ceiling equation forces $f=1$, again impossible.
Both index guards are restored.  The pinned local factorial lemma
now gives $f=K!$, so $f$ is even.  Its complete low-digit and
self-referential-radius proof is bundled; it uses these three calls,
not an interpretation of the remaining powers.

\paragraph{Orienting the t-pair with a different output.}
The base $x_{103}=1+(t+1)f$ is odd.  The t-seed norm and its
actual call-2 residue modulo two imply $x_{104}$ is odd, without
interpreting t.  Hence $t$ and $k_{122}=t+x_{104}+2$ have opposite
parity.  An even index excludes $|Y|=1$: the signed abscissa
half-period argument makes the comparison index odd, while its
ordinate congruence modulo four makes it even.  Thus its true index
is recovered, and a positive-output call orients its pair.

If t is even, use \emph{call 2}, whose output $x_{104}\ge0$ was
proved in \eqref{eq:direct-output104}; its base is at least five.
If $k_{122}$ is even, use \emph{call 9}, with output $x_{560}\ge0$.
The retained product of these two index slacks then makes the
other slack nonnegative.  The indexed t-pair is oriented by call 2
also when t is odd.  This change of supporting call avoids assuming
the missing sign of call 0's output $x_6$.

With that positive t-pair, the signed-output lemma at bases $t+1$
and $f$ recovers $x_6=(t+1)^t>0$ and $x_{106}=f^t>0$.
The remaining source rows give
\[
 x_{125}=x_{123}x_{107}x_{106}\ge0,\quad
 x_{125}=M+x_{126}x_{130},\quad
 x_{126}=M+x_{131}+1>M.
\]
Consequently $x_{130}\ge0$; its dependent signs follow.  The
entire 643-coordinate semantic prefix is now natural.

\paragraph{The second factorial's shifted index.}
It remains to recover the deleted predicate at 825.  With t exact,
actual calls 4 and 8 give
\[
 N_1=x_{113}=(1+f_1t^2)^t\ge1,\qquad x_{547}=b^t.
\]
The second factorial rows give
$C_1=x_{114}+x_{115}+1\ge1$ and $f_1C_1=N_1+x_{115}$,
so $f_1\ge1$.  Its shifted seed has index $t+f_1t^2$.
The large-ordinate case recovers its true positive pair as above.
In the small case $|Y|=1$, the signed residue analysis forces
$x_{546}=b+1$.  Its exact digit equation is
\[
 x_{546}=b^t(C_1+ba_1)+u_1,\qquad
 C_1\ge1,\quad a_1\ge0,\quad0\le u_1<b^t.
\]
For $t\ge2$ the right side is at least $b^t>b+1$.
For $t=1$, division forces $C_1=1$, whereas
$f_1C_1=f_1+1+x_{115}>f_1$.  Both cases are impossible.
Thus the omitted shifted-index slack is strictly positive.

\subsection{Completeness requires a fresh private frame}

All 75 parent predicate values have now been recovered.  The unchanged
Pell and small-residue arguments give every true power value; they
use only the dominating bounds, not the literal old linear frame.
Nevertheless, adjoining guard triples alone does \emph{not} produce
an old source zero: its original frame equality usually fails.

Hold the natural semantic prefix fixed and reconstruct the private
Pell data.  Choose either its old linear frame $W_0$ or the new squared
one, and set
\[
 A_0=\chi_{W_0}(W_0+1),\qquad
 Z_0=\psi_{W_0}(W_0+1)/W_0.
\]
The recurrence modulo $W_0$ proves the division is exact.  Choose
all main pairs at their actual indices.  In the finite quadratic
ring modulo $L=4\prod_iY_i^2$, the unit
$A_0+\sqrt{A_0^2-1}$ has finite order.  This supplies a positive
common Pell solution with $U\equiv1$, $V\equiv0\pmod L$.
The norm gives $\gcd(U,V)=1$, and $U$ is odd.  CRT supplies
$B\equiv A_0\pmod U$, $B\equiv1\pmod{4V}$, with $B>A_0$.
Take the comparison pairs at their true indices and use their
integral progression quotients.  The product quotient is
$V/(\prod_iY_i^2)-1\ge0$.

For each power call, its integral quotient follows from
\[
 H_0=1,\quad H_{j+1}=nH_j+T_n\psi_j(A_0),\quad
 H_j=\chi_j(A_0)-(A_0-n)\psi_j(A_0).
\]
Reevaluate the derived t+1 pair, discriminant, moduli and private
arithmetic gates.  Every one of the 109 private nongate residuals
is the sole nongate, with coefficient $\pm1$, in an exact row
identity for a target just satisfied.  All other terms are unit
gate equations.  The independent audit checks this entire coverage,
including the 474 private topological gate definitions.  Fresh
three-square witnesses satisfy every private guard.

Using the old linear frame constructs an old 75-predicate zero from
any new zero.  Conversely, use the old source's true powers and the
same completion with the squared frame, omitting the twelve predicates.
This constructs a new zero.  The maps establish existence equivalence;
they are not asserted to be inverse on private witnesses.

\subsection{Exact algebraic verification}

The final projections contain respectively 872 and 866 monic
eliminations.  Each deleted coordinate has a unique polynomial
integer restoration, every surviving residual stays quadratic,
and no unconstrained variable is discarded.  Reversing the sequence
gives a bijection of complete integer zero sets, including guard rows.
An independently authored checker verifies the source changes,
604-sign initial dependency boundary, all staged interfaces and
all 109 private completion targets.  A separate checker imports no
optimizer, replays every signed restoration and collects every
quartic coefficient.  Both normal and optimized Python modes pass,
and changed source rows, guards, inverses and coefficients are rejected.

These finite checks establish the literal polynomials and their
algebraic transformations.  The unbounded arguments are the Pell,
integer-order, local factorial and finite CRT proofs given above and
in the bundled pinned sources.  They do not prove Con(ZFC) itself.
The next concrete research questions are to minimize the outputs
needed to orient all seed pairs, and to optimize the signed frame
together with degree-preserving monic projection: a square is free
in the raw coordinate count but may prevent a later elimination.


\section{Verification, integration, and limitations}

The mathematical specification and reproducibility package are complementary.
The ten-positive polynomial has a complete closed coefficient recipe in
this report and a literal arithmetic graph in \file{ten_positive/equation.json}.
The direct quartics have complete residual systems, complete collected
monomial tables and all integer reconstruction maps in \file{direct/}.
No unknown coefficient or program index remains to be chosen later.

The checks performed for the ten-positive result include:
\begin{itemize}
\item the entire retained/changed boundary against the 99-gate parent;
\item exact new front-register identities and every coefficient of the
four-conjugate norm and double-center identity;
\item degree bounds for every centered summand, the complete nonzero leading
form, and its explicit degree-278 coefficient;
\item all four coefficient payloads reconstructed from the 436-character
seed, including the 943 tiles and 66,062 finite native convolution checks;
\item separate independent proof review of root extraction, dependency order,
finite-period witnesses and the auxiliary growth needed for $n>0$.
\end{itemize}

The direct route additionally checks every original-row identity used for
the changed bound and recovered predicates, every monic integer inverse,
every residual, and every coefficient of both expanded quartics. The normal
and optimized Python replays preserve the checks: essential verifications
do not disappear under \texttt{python -O}. Bounded examples and modular
samples are diagnostic controls, not substitutes for the unbounded proofs.

The archive includes a top-level \file{verify.py}, a source ledger, licenses
and deterministic build commands. The original repository is not required
to evaluate the shipped polynomials or run the standalone checks. The source
pin remains useful for reviewing the inherited semantic calibrations.

There are two principal limits. First, the ten-positive route still uses the
native K215 source calibration; its smaller arithmetic does not make it an
independent foundational encoding. The direct quartic route supplies that
different encoding. Second, neither new ordinary equivalence is presented
as a new end-to-end proof-assistant theorem. Exact algebra checks certify
the finite artifacts and support transcription accuracy; they do not
themselves certify all source semantics or prove $\Con(\mathrm{ZFC})$.

\section{Further research questions}

\begin{enumerate}
\item \textbf{Can the signed-output frame be generalized without creating
expensive quadratic dependencies?} Squaring an argument removes its sign
premise but may prevent an earlier monic quadratic substitution. The actual
variable count must be measured after projection. A joint choice of frame
terms and elimination order may improve the 576-variable profile.
\item \textbf{Can the two-root adjacent certificate eliminate another
positive unknown?} The present reduction compresses square conditions, not
the free source coordinates or the auxiliary index witness. A further
reduction must supply complete source reconstruction and canonical
completeness; a shorter norm formula alone is insufficient.
\item \textbf{Can better rational centers lower the degree again?} The
upper center $P=AI_0D$ has a degree-46 defect. A closer center may reduce that
defect, but clearing a new denominator can increase degree and interfere
with the projection's coprimality. Both effects must be counted together.
\item \textbf{Can saturation transfer to the signed low-variable branch?}
An unrestricted-integer witness absorbs changes to the center easily, but
some signed parents require an \emph{odd} strong multiplier. A finite-period
argument guarantees divisibility, not oddness. For example, at Pell base 2,
the least index with $7\mid\psi_2(m)$ is 4; all such indices are multiples
of 4 and their ordinates are even. An odd-index divisibility theorem is a
real additional obligation, not an automatic corollary of this packet.
\item \textbf{Can the ordinary source calibrations be internalized with
small certificates?} The local root extraction and finite-period lemmas
are natural targets once formalization resumes. The final internal-ZFC
statement must still quantify over all internally coded witnesses.
\end{enumerate}

\appendix
\section{A closed recipe for every coefficient}
\label{app:constants}

This appendix retains the entire coherent BFS-renamed compact44 recipe from
the pinned one-page source \cite{bfspage}. Its $K$ is the main equation's
$K_0$. All summation ranges, indices, sets, ranks, quotients and exponents
here are fixed or bound in the coefficient definitions; none is an extra
Diophantine unknown. The recipe was reconstructed and compared with every
raw payload byte. It is attributed as inherited data, not as a new machine
construction. In particular, the initializer uses the decoded first target.

\subsection{The four fixed integers}
Every index, set, sum, product and floor below defines a constant, not
another unknown. Put
\[
\rho=2^{67},\quad B=\rho^{8807},\quad\beta=B-1,\quad
G_m(x)=\sum_{j=0}^{m-1}x^j.
\]
Define
\[
\begin{split}
\mathcal E={}&\rg0{942}\cup\rg{983}{992}\cup\rg{1003}{1011}\\
&\cup\rg{1013}{1031}\cup\{1063,1064,1095,2195,3296\}.
\end{split}
\]
Then
\begin{align*}
M_C={}&B-1-\sum_{e\in\mathcal E\setminus\{1\}}\rho^e-2\rho^{1063},\\
M_F={}&B-1+4+(2^{51}-2)(\rho^{4393}+\rho^{5495})\\*
&+\rho^{1013}G_{10}(\rho)+\rho^{1065}G_9(\rho)
  +\rho^{2195}+\rho^{3296}\\*
&+(2^{65}-1)\rho^{8792}G_{13}(\rho)+(2^{19}-1)\rho^{8805},\\
K={}&\rho^{30}+\rho^{62}+\rho^{1100}+\rho^{2201}+\rho^{8849}\\*
&+(\rho^{4393}+\rho^{5495})\Phi(1/\rho),
\end{align*}
where
\[
\begin{split}
\Phi(x)={}&(1+2^{51})G_{943}(x)
 -2^{10}\sum_{t\in\mathcal T}v(t)x^{\varrho(t)}\\
&+\sum_{(b,e,m)\in\mathcal C}2^b x^e G_m(2x),\\
\mathcal C={}&\{(10,1023,8),(18,1013,9),(27,1003,8),\\
&(35,983,9),(44,1031,1),(45,1022,1),(46,1011,1),\\
&(47,992,1),(48,1095,1),(49,2195,1),(50,3296,1)\}.
\end{split}
\]
The rational intermediate in K cancels exactly; all four results are positive integers.

\subsection{The 943 fixed tiles and their ranks}
Use the actions $u_j,v_j,d_j$ decoded below. Let
\[
\begin{gathered}
h(q,b)=2(q+1)+b,\quad
\mathcal R=\{u_j:d_j=1,\ 0\le j<430\},\\
\mathcal L=\{u_j:d_j=0,\ 0\le j<430\},\\
r_q=6+|\mathcal R\cap\rg0{q-1}|,\quad
l_q=205+|\mathcal L\cap\rg0{q-1}|.
\end{gathered}
\]
An interval ending at $-1$ is empty. The following union is $\mathcal T$;
the index ranges in each row are fixed:
\[
\setlength{\arraycolsep}{1.4pt}\begin{array}{c|l}
\text{range}&\text{tile }(w,n,e,s)\\\hline
j=2q+b<430&((1-d_j)l_{u_j},h(q,b),d_jr_{u_j},v_j)\\
b=0,1&(0,h(215,b),0,h(215,b)),\ (0,b,0,b)\\
q\in\mathcal R,\ b=0,1&(r_q,b,0,h(q,b))\\
q\in\mathcal L,\ b=0,1&(0,b,l_q,h(q,b))\\
t\in\{0,1,432,433\}&(0,t,0,438)\\
&(2,434,2,0),\ (2,435,3,1)\\
&(3,436,4,h(u_0,0)),\ (4,436,4,0),\ (5,0,0,0)\\
t=434,435,436&(1,438,1,t)\\
a\in\{0,1,4\},\ b=0,1,2&(a,437,b,437)
\end{array}
\]
Put $s_*=(2,435,3,1)$, $f_*=(5,0,0,0)$, and $\mathcal T'=\mathcal T\setminus\{s_*,f_*\}$.
For $t=(w,n,e,s)$ set
\[
\begin{gathered}
\kappa(t)=2^{30}w+2^{20}n+2^{10}e+s,\\
v(t)=w+2^8n+2^{17}e+2^{25}s+127\,2^{34},\\
\varrho(s_*)=0,\quad\varrho(f_*)=1,\\
\varrho(t)=2+\sum_{z\in\mathcal T'}
 \fl{\frac{2^{40}+\kappa(t)-\kappa(z)-1}{2^{40}}}
 \quad(t\in\mathcal T').
\end{gathered}
\]

\subsection{A closed decoder for the 430 actions}
Put $D_b(A,k)=\fl{A/b^k}\bmod b$, where
$a\bmod b=a-b\fl{a/b}$. For $0\le q<215$, $b=0,1$, $j=2q+b$:
\[
\begin{gathered}
N=\mathcal A\bmod2^{636},\quad
T_0=\fl{\mathcal A/2^{636}}\bmod2^{1675},\\
E_0=\fl{\mathcal A/2^{2311}},\\
\nu_j=D_2(N,j),\quad \sigma_j=\sum_{k=0}^{j-1}\nu_k,\\
m_q=D_2(N,430+q),\quad \tau_q=\sum_{k=0}^{q-1}m_k,\\
F_q=14(1-m_q)+m_qD_{16}(E_0,\tau_q),\\
u_j=\nu_j(1+\sigma_j)+(1-\nu_j)D_{216}(T_0,j-\sigma_j),\\
v_j=D_2(F_q,2b),\quad d_j=D_2(F_q,2b+1).
\end{gathered}
\]
Empty sums are zero. Here $(v_0,d_0,u_0)=(1,1,1)$; the constructor tile
uses this actual first target, including its state renaming.

\subsection{The complete seed numeral}
The alphabet, in order of digit values $0,\ldots,61$, is
\begin{center}\small\ttfamily
0123456789ABCDEFGHIJKLMNOPQRSTUVWXYZabcdefghijklmnopqrstuvwxyz
\end{center}
Concatenate the following 436 characters as digits $e_0,\ldots,e_{435}$;
put $\mathcal A=\sum_{j=0}^{435}e_j62^{435-j}$.\par
\begingroup\small\ttfamily
BxEnhvy6NipzkWmROo0GsA3uK3twWLxscLq20UTumy5Cdu461LwD6g7RDNzF8Y\par
iYDFTpB8ehAypd7d6OsOL4APamnJeNXQypUezCmW9n4oXP6XfgmiWuAWZ8ZJ7U\par
OP7gu4kVdlGZqZHULQMar1dXNeJW4i7oNVI8usp8BlovjjZbwPCovZI328fRUq\par
4q92R81s5GjWQaQfGhkq9UYMK3rRMaBmsOIBkJqc2cj2VRIZ6mjzkBIOkW9SjZ\par
FzeA6J5u4j3sBL3SdgwXzM9mHvQivu8EJ1TBTpRhHDqZIbmYgwQ8VqTAFotl6F\par
We1YC554GKrVHsA8oUGKSA6F6k5KREBSD0Tgd7bFALY84LLbpw574BQyvdoHrl\par
0Oe35uwDf15DpT6VJWZBjSbHPpPAINYTte0eUxnTFTKBlGj8Am6sb1gLldh8WZ\par
8P\par
\endgroup


\section{File formats and reproduction}

An arithmetic graph row \texttt{[name,op,left,right]} defines one new
register using one of \texttt{+,-,*}; an operand is either an earlier name or
an integer literal. The output register is \texttt{N}. Coefficients are
unsigned little-endian integers stored in deterministic compressed payloads;
the closed recipe above also constructs them without reading those payloads.

The direct polynomial files use explicit sparse monomial lists. A monomial's
variable indices and exponents, together with its integer coefficient,
determine it completely. The residual systems are convenient for proof and
evaluation, while the expanded files permit independent coefficient-level
comparison. The direct package README specifies the exact JSON fields.

From the unpacked archive directory:
\begin{verbatim}
python3 -B verify.py
python3 -B -O verify.py
python3 -B ten_positive/build.py --check
python3 -B ten_positive/rebuild_coefficients.py
pdflatex -interaction=nonstopmode -halt-on-error report.tex
pdflatex -interaction=nonstopmode -halt-on-error report.tex
\end{verbatim}
The first two commands replay the shipped finite certificates. The report
does not require network access or a proof-assistant installation to build.
The full source ledgers distinguish copied upstream results from new files.

\begin{thebibliography}{9}
\bibitem{zerooffset}
ProveIt contributors, \emph{Ten strictly positive integers, 99 operations,
degree 368}, \file{Computability/DiophantineConsistency/candidates/ten-positive-zero-offset/},
at the fixed snapshot. \url{https://github.com/VladimirReshetnikov/ProveIt/tree/84e16e43092614bae262ca59651cb9d77e31a255/Computability/DiophantineConsistency/candidates/ten-positive-zero-offset}.
\bibitem{source61}
ProveIt research package 61, \emph{Congruence carriers},
\file{con-zfc-diophantine-forms/61-congruence-carriers/report/congruence_carriers.tex},
especially the source theorem, normalized rank, signed auxiliary index and
positive auxiliary reconstruction. The complete source is copied in
\file{ten_positive/upstream/SOURCE61.tex}.
\bibitem{bfspage}
ProveIt contributors, \emph{A complete ten-variable consistency equation},
\file{candidates/one-page-short10/}, and the inherited complete BFS source
transfer. \url{https://github.com/VladimirReshetnikov/ProveIt/tree/84e16e43092614bae262ca59651cb9d77e31a255/Computability/DiophantineConsistency/candidates/one-page-short10}.
\bibitem{ninesaturation}
ProveIt contributors, \emph{Tightened nine-positive main-root folding},
\file{research/nine-folded-saturation-tight/PROOF.md}, at the fixed snapshot.
The source proof is copied as \file{ten_positive/upstream/NINE_SATURATION_PRIOR.md}.
\bibitem{output75}
ProveIt contributors, \emph{Fresh 75-predicate integer quartic profiles},
\file{research/factorial-radius-output-sign-profile/}, at the fixed snapshot.
\url{https://github.com/VladimirReshetnikov/ProveIt/tree/84e16e43092614bae262ca59651cb9d77e31a255/Computability/DiophantineConsistency/research/factorial-radius-output-sign-profile}.
\bibitem{directproofs}
ProveIt contributors, direct C222 source, shared powers and shared index
certificates; signed comparison-base, paired-index, factorial-radius and
output-sign reductions. Exact source paths and hashes are included in
\file{direct/source/} and the direct source ledger.
\bibitem{sun2026}
Zhi-Wei Sun, \emph{Ten unknowns for Hilbert's tenth problem over the integers},
arXiv:2609.01594v2, 2 September 2026.
\url{https://arxiv.org/abs/2609.01594}.
\end{thebibliography}
\end{document}
